\subsection{Hölder's inequality}

In this No., p and q will denote two real numbers such that \(1 \leqslant p \leqslant +\infty\) , \(1 \leqslant q \leqslant +\infty\) , bound by the relation \(1/p + 1/q = 1\) ; thus \(q = p/(p - 1)\) if \(1 < p < +\infty\) , \(q = +\infty\) if p = 1, and q = 1 if \(p = +\infty\) ; p and q will be called conjugate exponents. Note that the relation \(1 \leqslant p \leqslant 2\) is equivalent to \(2 \leqslant q \leqslant +\infty\) ; p = q only when p and q are equal to 2.

THEOREM 2 (Hölder's inequality). --- Let \(f\) and \(g\) be two numerical functions that are finite almost everywhere and are such that \(f\) is equal almost everywhere to a function in \(\mathcal{L}^p\) and \(g\) to a function in \(\mathcal{L}^q\) . Then, the function \(fg\) (defined almost everywhere) is integrable, and

\[
\mathrm{N} _ {1} (f g) \leqslant \mathrm{N} _ {p} (f) \mathrm{N} _ {q} (g).\tag{4}
\]

Let \(f_{1}\) (resp. \(g_{1}\) ) be a function in \(\mathcal{L}^{p}\) (resp. \(\mathcal{L}^{q}\) ) to which \(f\) (resp. \(g\) ) is equal almost everywhere; \(fg\) is equal almost everywhere to the function \(f_{1}g_{1}\) , which is everywhere defined and finite, and which is measurable, being the product of two measurable functions (§5, No.~3, Cor. 5 of Th. 1). If \(1 < p < +\infty\) , Hölder's inequality for the upper integral (Ch. I, No.~3, Prop. 4) yields the inequality (4), and the relation \(N_{1}(fg) < +\infty\) then shows that \(fg\) is integrable (§5, No.~6, Th. 5). If \(p = 1\) , \(q = +\infty\) , the inequality (4) and the fact that \(fg\) is integrable are immediate consequences of the inequality of the mean (No.~2, Prop. 1); thus the theorem is proved in all cases.

COROLLARY 1. --- Let F, G, H be three Banach spaces, and let \((\mathbf{u}, \mathbf{v}) \mapsto \Phi(\mathbf{u}, \mathbf{v})\) be a continuous bilinear mapping of \(F \times G\) into H such that \(|\Phi(\mathbf{u}, \mathbf{v})| \leqslant |\mathbf{u}| \cdot |\mathbf{v}|\) . If \(f \in L_{F}^{p}\) and \(g \in L_{G}^{q}\) , then the function \(\Phi(\mathbf{f}, \mathbf{g})\) is integrable and

\[
\left| \int \Phi (\mathbf {f}, \mathbf {g}) d \mu \right| \leqslant \int | \Phi (\mathbf {f}, \mathbf {g}) | d | \mu | \leqslant \mathrm{N} _ {p} (\mathbf {f}) \mathrm{N} _ {q} (\mathbf {g}).\tag{5}
\]

For, \(\Phi(\mathbf{f},\mathbf{g})\) is measurable (§5, No.~3, Cor. 5 of Th. 1); since \(|\Phi(\mathbf{f},\mathbf{g})| \leqslant |\mathbf{f}| \cdot |\mathbf{g}|\) , the corollary follows from Th. 2 and the integrability criterion of §5, No.~6, Th. 5.

Two special cases of Cor. 1 are important in applications:

COROLLARY 2. --- Let F be a real (resp. complex) Banach space, F' its strong dual (TVS, III, §3, No.~1), and let \((\mathbf{z}, \mathbf{z}') \mapsto \langle \mathbf{z}, \mathbf{z}' \rangle\) be the canonical bilinear form on \(F \times F'\) . If \(f \in L_{F}^{p}\) and \(g \in L_{F'}^{q}\) , then the real (resp. complex) function \(\langle f, g \rangle\) is integrable and

\[
\left| \int \langle \mathbf {f}, \mathbf {g} \rangle d \mu \right| \leqslant \int | \langle \mathbf {f}, \mathbf {g} \rangle | d | \mu | \leqslant N _ {p} (\mathbf {f}) N _ {q} (\mathbf {g}).\tag{6}
\]

For, \(|\langle \mathbf{z},\mathbf{z}^{\prime}\rangle |\leqslant |\mathbf{z}|\cdot |\mathbf{z}^{\prime}|\)

When F is a real or complex Hilbert space, one knows that it can be canonically identified with its dual \(F'\) (TVS, V, §1, No.~7). Since the space \(L_{F}^{2}\) is complete, we have the following result:

COROLLARY 3. --- Let \(\mu\) be a positive measure on X, F a real (resp. complex) Hilbert space. On the space \(L_{F}^{2}\) , the symmetric (resp. Hermitian) form

\[
(\widetilde {\mathbf {f}}, \widetilde {\mathbf {g}}) \mapsto \int \langle \mathbf {f}, \mathbf {g} \rangle d \mu
\]

defines a Hilbert space structure, for which the norm is equal to \(\| \widetilde{\mathbf{f}}\| _2\)

COROLLARY 4. --- Let F be a Banach space, f a function in \(\mathcal{L}_{\mathrm{F}}^{p}\) , and g a numerical function belonging to \(\mathcal{L}^q\) ; then, the function fg is integrable and

\[
\left| \int \mathbf {f} g d \mu \right| \leqslant \int | \mathbf {f} g | d | \mu | \leqslant \mathrm{N} _ {p} (\mathbf {f}) \mathrm{N} _ {q} (g).\tag{7}
\]

COROLLARY 5. --- Let \(f_1, f_2, \ldots, f_n\) be \(n\) positive integrable functions, and let \(\alpha_1, \alpha_2, \ldots, \alpha_n\) be \(n\) numbers \(> 0\) such that \(\sum_{i=1}^{n} \alpha_i = 1\) ; under these conditions, the function \(f_1^{\alpha_1} f_2^{\alpha_2} \cdots f_n^{\alpha_n}\) is integrable and

\[
\begin{array}{l} \int f _ {1} ^ {\alpha_ {1}} f _ {2} ^ {\alpha_ {2}} \dots f _ {n} ^ {\alpha_ {n}} d | \mu | \leqslant \\ \left(\int f _ {1} d | \mu |\right) ^ {\alpha_ {1}} \left(\int f _ {2} d | \mu |\right) ^ {\alpha_ {2}} \dots \left(\int f _ {n} d | \mu |\right) ^ {\alpha_ {n}}. \end{array}\tag{8}
\]

For, the product \(f_{1}^{\alpha_{1}}f_{2}^{\alpha_{2}}\cdots f_{n}^{\alpha_{n}}\) is measurable, being the product of measurable functions (§5, No.~3, Th. 1 and its Cor. 5); since the inequality (8) is true for upper integrals (Ch. I, No.~2, Cor. of Prop. 2), the function \(f_{1}^{\alpha_{1}}f_{2}^{\alpha_{2}}\cdots f_{n}^{\alpha_{n}}\) is integrable (§5, No.~6, Th. 5), whence the corollary.

Cor. 2 of Th. 2 is sharpened by the following proposition:

PROPOSITION 3. --- Let \(\mu\) be a positive measure on X, F a real or complex Banach space, F' its strong dual, and \((\mathbf{z},\mathbf{z}')\mapsto\langle\mathbf{z},\mathbf{z}'\rangle\) the canonical bilinear form on F × F'.

\(1^{\circ}\) For every function \(\mathbf{f} \in \mathcal{L}_{\mathrm{F}}^{p}\) ( \(1 \leqslant p \leqslant +\infty\) ),

\[
\mathrm{N} _ {p} (\mathbf {f}) = \sup \left| \int \langle \mathbf {f}, \mathbf {g} \rangle d \mu \right|\tag{9}
\]

as \(\mathbf{g}\) runs over the set of functions in \(\mathcal{L}_{\mathrm{F}'}^q\) , such that \(\mathrm{N}_q(\mathbf{g}) \leqslant 1\) .

\(2^{\circ}\) For every function \(\mathbf{g} \in \mathcal{L}_{\mathrm{F}}^{q}\) ( \(1 \leqslant q \leqslant +\infty\) ),

\[
\mathrm{N} _ {q} (\mathbf {g}) = \sup \left| \int \langle \mathbf {f}, \mathbf {g} \rangle d \mu \right|\tag{10}
\]

as \(\mathbf{f}\) runs over the set of functions in \(\mathcal{L}_{\mathrm{F}}^{p}\) such that \(\mathrm{N}_p(\mathbf{f})\leqslant 1\)

Let us first prove the relation (9); we distinguish between two cases.

\begin{enumerate}
\def\labelenumi{(\roman{enumi})}
\tightlist
\item
  \(1 \leqslant p < +\infty\) . The relation (9) being trivial when \(\mathrm{N}_p(\mathbf{f}) = 0\) (because \(\mathbf{f}\) and \(\langle \mathbf{f}, \mathbf{g} \rangle\) are then negligible), we can always suppose, on multiplying \(\mathbf{f}\) by a scalar, that \(\mathrm{N}_p(\mathbf{f}) = 1\) . Suppose first that \(\mathbf{f}\) is an integrable step function, \(\mathbf{f} = \sum_{k=1}^{n} \dot{\mathbf{a}}_k \varphi_{\mathrm{A}_k}\) , where the \(\mathrm{A}_k\) are pairwise disjoint (§4, No.~9, Lemma). Thus \(\sum_{k=1}^{n} |\mathbf{a}_k|^p \mu(\mathrm{A}_k) = 1\) by hypothesis. For every \(\varepsilon > 0\) , there exists (for every index \(k\) ) a vector \(\mathbf{a}_k' \in \mathrm{F}'\) such that \(|\mathbf{a}_k'|^q = |\mathbf{a}_k|^p\) if \(p > 1\) (resp. \(|\mathbf{a}_k'| = 1\) if \(p = 1\) ) and \(\langle \mathbf{a}_k, \mathbf{a}_k' \rangle \geqslant (1 - \varepsilon) |\mathbf{a}_k| \cdot |\mathbf{a}_k'|\) (TVS, IV, §1, No.~3, Prop. 8). Setting \(\mathbf{g} = \sum_{k=1}^{n} \mathbf{a}_{k}^{\prime} \varphi_{\mathrm{A}_{k}}\) , we have \(\sum_{k=1}^{n} |\mathbf{a}_{k}^{\prime}|^{q} \mu(\mathrm{A}_{k}) = 1\) if \(p > 1\) (resp. \(\sup_{1 \leqslant k \leqslant n} |\mathbf{a}_{k}^{\prime}| = 1\) if \(p = 1\) ), thus \(\mathrm{N}_{q}(\mathbf{g}) = 1\) ; on the other hand,
\end{enumerate}

\[
\int \langle \mathbf {f}, \mathbf {g} \rangle d \mu = \sum_ {k = 1} ^ {n} \left\langle \mathbf {a} _ {k}, \mathbf {a} _ {k} ^ {\prime} \right\rangle \mu \left(\mathrm{A} _ {k}\right) \geqslant (1 - \varepsilon) \sum_ {k = 1} ^ {n} | \mathbf {a} _ {k} | \cdot | \mathbf {a} _ {k} ^ {\prime} | \mu (\mathrm{A} _ {k})
\]

and, since \(|\mathbf{a}_k'| = |\mathbf{a}_k|^{p / q} = |\mathbf{a}_k|^{p - 1}\) if \(p > 1\) (resp. \(|\mathbf{a}_k'| = 1\) if \(p = 1\) ), we have

\[
\int \langle \mathbf {f}, \mathbf {g} \rangle d \mu \geqslant (1 - \varepsilon) \sum_ {k = 1} ^ {n} | \mathbf {a} _ {k} | ^ {p} \mu (\mathrm{A} _ {k}) = (1 - \varepsilon) \mathrm{N} _ {p} (\mathbf {f}) = 1 - \varepsilon ,
\]

which proves the relation (9) in this case.

Let us pass to the case that \(\mathbf{f}\) is any element of \(\mathcal{L}_{\mathrm{F}}^{p}\) such that \(\mathrm{N}_p(\mathbf{f}) = 1\) . For every \(\varepsilon > 0\) , there exists a step function \(\mathbf{f}_1 \in \mathcal{L}_{\mathrm{F}}^p\) such that \(\mathrm{N}_p(\mathbf{f} - \mathbf{f}_1) \leqslant \varepsilon\) (§4, No.~10, Cor. 1 of Prop. 19). By what we have just seen, there exists a function \(\mathbf{g} \in \mathcal{L}_{\mathrm{F}}^q\) , such that \(\mathrm{N}_q(\mathbf{g}) = 1\) and

\[
\int \langle \mathbf {f} _ {1}, \mathbf {g} \rangle d \mu \geqslant \mathrm{N} _ {p} (\mathbf {f} _ {1}) - \varepsilon \geqslant 1 - 2 \varepsilon .
\]

Now,

\[
\int \left\langle \mathbf {f}, \mathbf {g} \right\rangle d \mu = \int \left\langle \mathbf {f} _ {1}, \mathbf {g} \right\rangle d \mu + \int \left\langle \mathbf {f} - \mathbf {f} _ {1}, \mathbf {g} \right\rangle d \mu
\]

and, by (6),

\[
\left| \int \langle \mathbf {f} - \mathbf {f} _ {1}, \mathbf {g} \rangle d \mu \right| \leqslant \mathrm{N} _ {p} (\mathbf {f} - \mathbf {f} _ {1}) \mathrm{N} _ {q} (\mathbf {g}),
\]

whence

\[
\left| \int \langle \mathbf {f}, \mathbf {g} \rangle d \mu \right| \geqslant 1 - 3 \varepsilon ,
\]

which proves (9).

\begin{enumerate}
\def\labelenumi{(\roman{enumi})}
\setcounter{enumi}{1}
\tightlist
\item
  \(p = +\infty\) . We may again restrict ourselves to the case that \(\mathrm{N}_{\infty}(\mathbf{f}) > 0\) . Let \(\alpha\) be any number such that \(0 < \alpha < \mathrm{N}_{\infty}(\mathbf{f})\) ; by hypothesis, the set of \(x \in X\) such that \(|\mathbf{f}(x)| > \alpha\) is measurable and is not locally negligible, therefore it contains a compact set K of measure \textgreater{} 0. Since f is measurable, there exists a compact set \(K_{1} \subset K\) of measure \textgreater{} 0, such that the restriction of f to \(K_{1}\) is continuous. It follows that, for every \(\varepsilon > 0\) , there exists a partition of \(K_{1}\) into a finite number of integrable sets, in each of which the oscillation of f is \(\leqslant \varepsilon\) ; at least one of these sets A has measure \(>0\) . Let \(\mathbf{a}\) be one of the values of \(\mathbf{f}\) in \(\mathbf{A}\) ; then \(|\mathbf{a}| > \alpha\) and \(|\mathbf{f}(x) - \mathbf{a}| \leqslant \varepsilon\) for all \(x \in \mathbf{A}\) . There exists a vector \(\mathbf{a}' \in \mathbf{F}'\) such that \(|\mathbf{a}'| = 1\) and \(|\langle \mathbf{a}, \mathbf{a}' \rangle| \geqslant |\mathbf{a}| - \varepsilon\) ; the function \(\mathbf{g} = \varphi_{\mathrm{A}} \cdot \mathbf{a}' / \mu(\mathbf{A})\) is integrable and \(\mathrm{N}_1(\mathbf{g}) = 1\) ; on the other hand,
\end{enumerate}

\[
\int \langle \mathbf {f}, \mathbf {g} \rangle d \mu = \frac {1}{\mu (\mathrm{A})} \int \langle \mathbf {f}, \mathbf {a} ^ {\prime} \rangle \varphi_ {\mathrm{A}} d \mu .
\]

Now, one can write

\[
\int \langle {\bf f}, {\bf a} ^ {\prime} \rangle \varphi_ {\mathrm{A}} d \mu = \langle {\bf a}, {\bf a} ^ {\prime} \rangle \mu (\mathrm{A}) + \int \langle {\bf f} - {\bf a}, {\bf a} ^ {\prime} \rangle \varphi_ {\mathrm{A}} d \mu ,
\]

and since

\[
\left| \langle \mathbf {f} - \mathbf {a}, \mathbf {a} ^ {\prime} \rangle \varphi_ {\mathrm{A}} \right| \leqslant \varepsilon \varphi_ {\mathrm{A}},
\]

we see that

\[
\left| \int \langle \mathbf {f}, \mathbf {g} \rangle d \mu \right| \geqslant | \langle \mathbf {a}, \mathbf {a} ^ {\prime} \rangle | - \varepsilon \geqslant | \mathbf {a} | - 2 \varepsilon > \alpha - 2 \varepsilon ;
\]

since \(\varepsilon\) is arbitrary and \(\alpha\) is any number \(<\mathrm{N}_{\infty}(\mathbf{f})\) , the relation (9) is also verified in this case.

One argues in exactly the same manner to prove the relation (10), on considering separately the case \(1 \leqslant q < +\infty\) and the case \(q = +\infty\) , and using the fact that for every \(\mathbf{z}' \in \mathrm{F}'\) , \(|\mathbf{z}'| = \sup_{|\mathbf{z}| \leqslant 1} |\langle \mathbf{z}, \mathbf{z}' \rangle|\) by the definition of the norm in \(\mathrm{F}'\) .

Remarks. --- 1) Let \(\mathcal{E}\) be a dense linear subspace of \(\mathcal{L}_{\mathrm{F}'}^q\) ; then the formula (9) holds when \(\mathbf{g}\) runs over the intersection of \(\mathcal{E}\) with the set B of functions in \(\mathcal{L}_{\mathrm{F}'}^q\) such that \(\mathrm{N}_q(\mathbf{g}) \leqslant 1\) . For, it suffices to observe that the interior \(\overset{\circ}{\mathrm{B}}\) of B is dense in B and that \(\overset{\circ}{\mathrm{B}} \cap \mathcal{E}\) is dense in \(\overset{\circ}{\mathrm{B}}\) , since \(\overset{\circ}{\mathrm{B}}\) is open. This remark applies in particular to the set \(\mathcal{E} = \mathcal{K}(\mathrm{X};\mathrm{F}')\) of continuous functions with compact support (with values in \(\mathrm{F}'\) ) when \(1 \leqslant q < +\infty\) , that is, \(1 < p \leqslant +\infty\) . But in this case, the formula (9) is true as \(\mathbf{g}\) runs over \(\mathrm{B} \cap \mathcal{K}(\mathrm{X};\mathrm{F}')\) , even for \(p = 1\) . For, we may, as above, restrict ourselves to the case that \(\mathbf{f}\) is a step function. We saw then that if \(\mathrm{N}_1(\mathbf{f}) = 1\) , then for every \(\varepsilon > 0\) there exists a step function \(\mathbf{g} \in \mathcal{L}_{\mathrm{F}'}^{\infty}\) such that \(|\mathbf{g}(x)| \leqslant 1\) for all \(x \in \mathrm{X}\) and \(|\int \langle \mathbf{f},\mathbf{g}\rangle d\mu| \geqslant 1 - \varepsilon\) . There exists a finite number of pairwise disjoint compact sets \(\mathrm{K}_i\) such that \(\mathbf{g}\) has a constant value \(\mathbf{a}_i'\) on each \(\mathrm{K}_i\) and such that, if K is the union of the \(\mathrm{K}_i\) , then \(\int |\mathbf{f}| \varphi_{\mathbf{C} K} d\mu \leqslant \varepsilon\) . Let \(\mathrm{U}_i\) be a neighborhood of \(\mathrm{K}_i\) such that the sets \(\mathrm{U}_i\) are pairwise disjoint, and let \(h_i\) be a continuous mapping of X into {[}0,1{]} with support contained in \(\mathrm{U}_i\) and equal to 1 on \(\mathrm{K}_i\) . Setting \(\mathbf{h} = \sum \mathbf{a}_i' h_i\) , we have \(\mathbf{h}(x) = \mathbf{g}(x)\) on K and \(|\mathbf{h}(x)| \leqslant 1\) on X, therefore

\[
\int | \langle \mathbf {f}, \mathbf {h} \rangle | \varphi_ {\mathbf {C K}} d \mu \leqslant \varepsilon
\]

and consequently \(|\int \langle \mathbf{f},\mathbf{h}\rangle d\mu |\geqslant 1 - 3\varepsilon\) , which proves our assertion. Analogous remarks can be made for the formula (10).

\begin{enumerate}
\def\labelenumi{\arabic{enumi})}
\setcounter{enumi}{1}
\tightlist
\item
  Let \(\mu\) be a positive measure on \(X\) , \(f\) a measurable function \(\geqslant 0\) (finite or not) whose support is contained in a countable union of compact sets \(K_{n}\) . Then, for every \(p\) such that \(1 \leqslant p \leqslant +\infty\) ,
\end{enumerate}

\[
\mathrm{N} _ {p} (f) = \sup \int^ {*} | f g | d \mu ,\tag{11}
\]

as \(g\) runs over the set of functions in \(\mathcal{K}(\mathrm{X};\mathbf{R})\) such that \(\mathrm{N}_q(g)\leqslant 1\) . For, the formula (11) is a special case of (9) when \(\mathrm{N}_p(f) < +\infty\) , since \(f\) is then equivalent to a function in \(\mathcal{L}^p\) (§5, No.~6, Th. 5). If \(\mathrm{N}_p(f) = +\infty\) , for every integer \(n > 0\) set \(f_{n} = \inf (n,f\varphi_{\mathrm{K}_{n}})\) . Then

\[
\mathrm{N} _ {p} (f _ {n}) = \sup \int^ {*} | f _ {n} g | d \mu \leqslant \sup \int^ {*} | f g | d \mu ,
\]

whence, on passing to the limit (assuming, as we may, that the sequence \((\mathbf{K}_n)\) is increasing), we have \(\sup \int^{*}|fg| d\mu = +\infty\) (§1, No.~3, Th. 3).

COROLLARY. --- Let \(\mu\) be a positive measure on X, F a Banach space, \(F'\) its strong dual, and g any function in \(L_{F'}^{q}\) . Then, the linear form on \(L_{F}^{p}\) , deduced from the linear form \(f \mapsto \int \langle f, g \rangle d\mu\) on \(L_{F}^{p}\) by passage to the quotient, is continuous and has norm \(N_{q}(g)\) .

